\subsection*{Appendix A: GPU Kernels with Triton}

\paragraph{Exercise A.1.} $\lceil10^6/1024\rceil=977$ programs, covering
$977\cdot1024=1{,}000{,}448$ lanes, so 448 lanes are masked. \code{tl.arange} requires a
power-of-two extent: register layouts distribute a block evenly over threads, warps and
vector lanes, all of which come in powers of two.

\paragraph{Exercise A.2.} With \code{other=0.0} the padded lanes take part in both
reductions. The max is harmless (softmax is shift-invariant, though all-negative rows
then get a needlessly shifted max), but the sum gains $e^{0-m}$ per padded lane. For the
row $(1,2,3)$ with \code{BLOCK}$=4$: the true denominator is $e^{-2}+e^{-1}+1\approx1.503$,
but the kernel adds $e^{-3}\approx0.050$. Every output is then too small by a factor
$1.503/1.553$ and the stored row sums to $0.968$.

\paragraph{Exercise A.3.} Loop over the row in chunks of \code{BLOCK} with running
$(m,\ell)$ updated by online rescaling, then loop again computing
$e^{x-m}/\ell$ and storing it. That is two reads and one write of the row. A single
write is natural; a single \emph{read} is possible only if the row is cached in SMEM
between the passes, which the compiler cannot do for unbounded rows. For moderate
lengths, rely on L2 to serve the second read.

\paragraph{Exercise A.4.} $I=2B_MB_N/((B_M+B_N)b)$. In bf16: 32, 64 and 85 FLOP/B; in
fp8: 64, 128 and 171. The H100's ridge points are about 295 (bf16) and 591 (fp8), so none
of these tiles could run at peak if every load came from HBM. The answer changes because
panels of $A$ and $B$ are re-read by many programs: with good ordering most re-reads hit
in L2, whose bandwidth is several times HBM's. The relevant intensity is then
whole-problem FLOPs over HBM bytes, which is $n/3$ for a square GEMM.

\paragraph{Exercise A.5.} With row-major order the first $P$ programs cover
$\lceil P/C\rceil$ rows and $\min(P,C)$ columns of tiles, which touches that many $A$ and
$B$ panels. With group size $G$ ($G\le R$, $P/G\le C$) they cover $G$ rows and
$\lceil P/G\rceil$ columns. The total $G+P/G$ is minimised at $G=\sqrt P$, with value
$2\sqrt P$ instead of up to $P+1$. For $P=9$: $3+3=6$ panels instead of 8 for
$C=6$ (\cref{fig:triton:swizzle}).

\paragraph{Exercise A.6.} Each stage buffers one $A$ and one $B$ slab:
$s(B_MB_K+B_KB_N)\cdot2$ bytes, i.e.\ $32$\,KB per stage for $128\times128\times64$.
With 228\,KB, $s\le7$, leaving one program per SM. Typical choices are $s=3$ (96\,KB,
two programs per SM) or $s=4$ (128\,KB, one program per SM). Remember that the
epilogue may need SMEM for the output tile as well (it is usually aliased onto
pipeline buffers).

\paragraph{Exercise A.7.} Load the packed tile (\code{BK/2}$\times$\code{BN} bytes),
split nibbles with \code{\& 0xF} and \code{>> 4}, and interleave them back to a
\code{BK}$\times$\code{BN} tile of codes. When the group size equals \code{BK}, each $K$
iteration has exactly one scale and zero point per output column. Load those two
\code{BN}-vectors once per iteration. Either dequantize before the dot,
$w=(q-z)\cdot s$, or compute \code{tl.dot(x, q - z)} and multiply the
\code{BM}$\times$\code{BN} partial product by $s$ before adding it to \code{acc}. The
latter applies the scale to $B_MB_N$ values instead of $B_KB_N$. The scale cannot be
deferred to the epilogue as in W8A16, because it changes along $K$.

\paragraph{Exercise A.8.} $Q$ and $O$ cost $2\cdot H\,T\,d\cdot2$\,B$=134$\,MB. Each of
the $H\,T/B_M$ programs streams all of $K$ and $V$ ($2S\,d\cdot2$\,B$=4.2$\,MB per head),
so without any cache hits the total is $(T/B_M)\cdot134$\,MB: 17.2\,GB for $B_M=64$ and
8.6\,GB for $B_M=128$ (half of that with causal skipping). Naive attention writes and
re-reads a 4.3\,GB score matrix and a 4.3\,GB probability matrix, about 17\,GB in total.
So without L2 the tiled kernel would \emph{not} save HBM traffic! It wins because one
head's $K$ and $V$ (4.2\,MB) fit comfortably in the 50\,MB L2: programs for the same head
run concurrently and reuse them. The real HBM traffic is then close to the 268\,MB of
$Q,K,V,O$. This is also why the grid puts the query-block index first (adjacent
programs share a head).

\paragraph{Exercise A.9.} With $T=1$ there is exactly one query block per (batch, head),
so the grid is $(1,B\cdot H)$, often fewer programs than SMs, each streaming a long KV.
For split-KV, chunk $c$ returns its running max $m_c$, normaliser $\ell_c$ and
unnormalised accumulator $\vect a_c$. A second kernel computes $M=\max_c m_c$,
$L=\sum_c e^{m_c-M}\ell_c$ and $\vect o=\sum_c e^{m_c-M}\vect a_c/L$. This is exact by
\cref{prop:triton:online} applied to chunks instead of blocks. (Equivalently, store
$\vect a_c/\ell_c$ and the log-sum-exp $m_c+\log\ell_c$.)

\paragraph{Exercise A.10.} Interleave the columns of $W_g$ and $W_u$ so that every
output tile contains matching gate and up columns, for example the first $B_N/2$ columns
of the tile from $W_g$ and the next $B_N/2$ from the same feature range of $W_u$. The
epilogue then computes $\SiLU(g)\odot u$ and writes a $B_M\times B_N/2$ tile. The
constraint is that $B_N$ must cover pairs, so tiles see half as many output features
per program. Per token and layer, the unfused path writes $g$ and $u$
($2\cdot12288\cdot2$\,B$=49$\,KB), reads them back (49\,KB) and writes $h$ (24.6\,KB).
The fused path only writes $h$, saving about 98\,KB per token per layer, or
$\approx3.5$\,MB per token over 36 layers.

\paragraph{Exercise A.11.} From $r=(\frac1n\sum_kx_k^2+\epsilon)^{-1/2}$ we get
$\partial r/\partial x_j=-r^3x_j/n$. So
$\partial\hat x_i/\partial x_j=r\delta_{ij}-x_ir^3x_j/n=r(\delta_{ij}-\hat x_i\hat x_j/n)$,
and $\partial\mathcal L/\partial x_j=\sum_ig_i\,\partial\hat x_i/\partial x_j
=r\big(g_j-\hat x_j\frac1n\sum_ig_i\hat x_i\big)$ with $g_i=dy_i\,w_i$. Also
$\partial\mathcal L/\partial w_i=\sum_{\text{rows}}dy_i\hat x_i$. Recomputing $r$ is
safe because it is a deterministic function of $x$, which the backward kernel loads
anyway. It is cheaper because a row reduction on data already in registers is a few
FLOPs per element, while saving $r$ would cost an extra write in the forward, a read in
the backward, and activation memory in between.

\paragraph{Exercise A.12.} Our kernel layout is \code{(num\_blocks, block\_size, H, D)}.
Per block, one program loads $16\cdot128\cdot2$\,B$=4$\,KB of $K$ and as much of $V$:
8\,KB, in 256-byte rows strided by $H\cdot D$. (A head-major block layout makes this
fully contiguous.) With GQA, one program per (sequence, KV head) gives $2\cdot8=16$
programs, about 12\% of a 132-SM GPU, each streaming 32k tokens. Split each sequence
into 16--32 chunks of 1--2k tokens, giving 256--512 programs (2--4 per SM), and merge
with the reduction from Exercise~A.9.
